\chapter{Amulets across different social, temporal and regional contexts}


This chapter presents a comparative analysis of the different objectscapes established in the previous chapters. The data will be analyzed across three main elements: social, chronological and regional. The center of the analysis is the social groups, which will then be evaluated within different chronological phases and locations. The chronological division follows the same one established in previous chapters. However, in order to improve direct comparisons---as many elite and non-elite tombs were dated over a broad range of phases (e.g. 25th Dynasty-Middle Napatan, 25th Dynasty-Early Napatan, etc.)---, the data was divided between the earlier and later phases. The regional division is as follows, from north to south of the Sudanese territory: Lower Nubia (Mirgissa, Qustul and Debeira 176), North Upper Nubia (Missiminia), 4th Cataract (El-Kurru, Nuri and Sanam) and Meroe region (Meroe West and South). The analysis comprehends the amulet types, motifs and material composition. The broader types are the same as in previous chapters, while the motifs were divided into interpretative subcategories to support the analysis, which will be further detailed in their specific sections. The materials were primarily divided into faience (the most common) and the remaining according to their places of origin, such as Lower Nubia, Egyptian Eastern Desert, Asia, among others, with the objective of understanding trade networks. Although, it is difficult to ascertain whether the materials or finished amulets were imported, as available evidence comes only from the Sanam Treasury and temple workshop, where raw materials and finished amulets were found.


\section{Chronological variations}


This section focus on the chronological variations of amulets across royal, elite and non-elite assemblages. The main objective is to interpret these variations based on the historical context of each phase and the internal social heterogeneity within Nubia. As a result, we can establish specific objectscapes of each different context in relation to amulets.


\subsection{Amulet distribution}


First, it is important to understand the distribution of amulets, especially their frequency in burials. As shown in Figure \ref{fig:phase-distro}, amulets were not included in the majority of burials, with elite and non-elite contexts displaying the lowest distribution. Only in royal sites of the 25th Dynasty and Early Napatan phases more than half of the tombs contained at least one amulet---70\% and 52\%, respectively. In fact, royal contexts display the biggest percentage of tombs with amulets, with the exception of those from the pre-25th Dynasty (8\%). In contrast, less than 12\% of elite tombs from all phases contained any amulets, with the highest percentage stemming from Late Napatan sites. For the non-elite, distribution follows the royal pattern, with more amulets present in tombs dated to the 25th Dynasty and Early Napatan phases (24\% and 22\%), while fewer tombs from the remaining phases contained amulets, only 5\% from pre-25th Dynasty and 6\% from Middle Napatan.


\begin{figure}[H]
	\centering
	\includegraphics[width=1\textwidth]{phase-distro}
	\caption{Distribution (in percentage) of tombs containing amulets by social group and chronological phase.}
	\label{fig:phase-distro}
\end{figure}


\begin{figure}[H]
	\centering
	\includegraphics[width=1\textwidth]{intact-meroe}
	\caption{Comparison of total number of intact tombs and those containing amulets in Meroe cemeteries by social group and chronological phase.}
	\label{fig:intact-meroe}
\end{figure}


The distribution could have been affected by the level of looting the tombs have suffered, especially the royal pyramids. Therefore, it is important to look into the intact tombs, which were located only at the Meroe cemeteries, as noted in the reports by Reisner/Dunham. As mentioned in Chapter 2, the other sites were also disturbed, but the reports do not explicitly note if any tombs were intact.


Taking into account the intact tombs from Meroe, it appears that non-elite graves were not a priority for plunderers (Figure \ref{fig:intact-meroe}). This can be attributed to their lack of superstructures, which would make than less visible, and therefore less desirable. Although the number of intact tombs was high, amulets were less frequent. Less than 20\% of 25th Dynasty graves contained amulets, while for Early Napatan the percentage was 15\% and for Middle Napatan, 19\%. There were nine intact elite burials which included amulets. One was from the 25th Dynasty and out of the eight total intact tombs from the Early Napatan phase, there were only two examples where amulets were found (25\%). Considering these figures, there is reason to believe that amulets were not an integral part of the funerary equipment, especially for the elite and non-elite. It is possible that their inclusion was based on personal choice and, consequently, their distribution varies on an individual grave basis.


\subsection{Amulet types}


\begin{figure}[H]
	\centering
	\includegraphics[width=1\textwidth]{phase-types}
	\caption{Frequency (in percentage) of amulet categories by social group and chronological phase.}
	\label{fig:phase-types}
\end{figure}


Regarding the distribution of amulet types, certain patterns and variations can be observed across social groups (Figure \ref{fig:phase-types}). In royal contexts, symbols are the most popular category across all phases. They were mostly followed by deities, with the exception of the Middle Napatan phase, when animals came in second. Some categories are totally absent in a given phase. For instance, amulets of animals, objects and nature were not found in pre-25th Dynasty and Late Napatan burials, whereas objects were absent from Middle Napatan graves. The most diverse assemblages come from 25th Dynasty and Early Napatan contexts, encompassing almost all seven categories (amulets of humans were only found in elite and non-elite burials). In contrast, pre-25th Dynasty and Middle Napatan assemblages included three amulet types, while Late Napatan was the least diverse, with only amulets of symbols and deities.


Similarly, symbol amulets were the most common in elite tombs throughout the Napatan period, with deities and animals appearing as the second and third most frequent categories from the 25th Dynasty to Middle Napatan phases. Meanwhile, Late Napatan assemblages were the least diverse, containing only symbols and animals. Six of the seven amulet types were present in 25th Dynasty and Middle Napatan assemblages. The earlier phase did not contain any amulets of objects, while human figures were absent from the latter. Early Napatan amulets were the most diverse, with the total seven categories present.


Non-elite assemblages differ significantly from the others, especially during the pre-25th Dynasty and Middle Napatan phases, when deities were the most common amulet type. Additionally, the second most common type during the Middle Napatan phase were human figures and not symbols, as is the case in the pre-25th Dynasty. 25th Dynasty and Early Napatan amulets follow the same pattern as those from contemporary royal and elite assemblages, when the three most popular types were symbols, deities and animals, in this order. Both phases were also the most diverse, including all seven amulet categories. In contrast, Middle Napatan burials included six amulet types, while those from the pre-25th Dynasty had only three: symbols, deities and animals. Compared to elite burials, there is an increase in the number of human amulets in non-elite contexts from the 25th Dynasty period, peaking in the Middle Napatan, when it became the second most common type.


Chronologically, the variety of amulets types is lower during the pre-25th Dynasty for all groups and reaches its peak during the 25th Dynasty and Early Napatan phases, before dropping again in Middle and Late Napatan assemblages. However, differences between social groups exist. While symbols were the majority of amulets in royal and elite contexts, in non-elite assemblages this only happens during the 25th Dynasty and Early Napatan phases. This may be a result of the control over Egypt, as well as ongoing trade relations afterwards. Concerning specific types, the main differences lie in the absence of human amulets in royal burials, as well as a higher frequency of animal and deity motifs in elite and non-elite contexts.


\subsection{Amulet motifs}


As previously mentioned, in order to better analyze the distribution of motifs, they were divided into interpretative subcategories based on their broader types. This choice was made with the aim to explore patterns and differences between social groups, as well as examining the impact of changing relations with Egypt during the Napatan period.


Deity amulets were divided into two categories: those from the Egyptian pantheon and local adaptations and/or deities. The former are those that are Egyptian in origin, while local adaptations/deities are depictions that differ from the standard Egyptian design, although oftentimes they were paired with Egyptian symbols. For instance, among those are the group of winged goddesses and gods, a ram-headed dwarf, unidentified deities---such as those with lion, snake, hawk, double hawk, crocodile and ram heads---, and adaptations of the Isis nursing Horus motifs (such as the lion-headed goddess/Horus and Isis/Nubian queen).


Overall, Egyptian deities were the most widespread across all phases and groups, but local adaptation/deities varied significantly across groups and phases (Figure \ref{fig:phase-deities}). In royal contexts, during the pre-25th Dynasty, the only deity amulet was of an unidentified hawk-headed dwarf back-to-back with a Pataikos, which I have considered as a local adaptation. In the 25th Dynasty, amulets included more local designs, but they still accounted for only 8\% of the total number of amulets against 19\% of Egyptian deities. Conversely, no recognizable local motifs were found in the subsequent royal phases, although they did include varying numbers of Egyptian deities: 4\% in the Early Napatan phase, 2\% in Middle Napatan and 12\% in Late Napatan.


In elite graves, local designs occur solely during the Early and Middle Napatan phases, 1\% and 6\%, respectively. These included two unknown deities, one with a hawk head and another with a snake head. Egyptian deities totaled 28\% in the 25th Dynasty burials, 14\% in Early Napatan and 15\% in Middle Napatan. For the non-elite, the occurrence of these amulets was even rarer, less than 1\% in two phases, 25th Dynasty and Early Napatan. The motifs encompass the leonine goddess nursing Horus, deities with the heads of a hawk, lion, snake and crocodile, two winged lion-headed goddess and a winged anthropomorphic goddess. Although the absolute quantity of these amulets is higher than in other social groups, the total amount of non-elite amulets is so high that their percentage is very low.


The data indicates that local adaptation/deities were much less common than Egyptian style motifs, but were proportionally more frequent in royal contexts of the 25th Dynasty, and especially reserved to queens and minor wives. They were also concentrated between the 25th Dynasty and Middle Napatan phases, although with significant differences between groups. It is important to note that, although certain deities were originally Egyptian, they likely were common in Nubia during the Napatan period, as their presence in the region was attested since the New Kingdom occupation and some even from the Kerma period. This notion will be further explored in Chapter 7.


\begin{figure}[H]
	\centering
	\includegraphics[width=1\textwidth]{phase-deities}
	\caption{Proportion (in percentage) of amulets of Egyptian deities and local adaptations by social group and chronological phase, relative to the total number of amulets.}
	\label{fig:phase-deities}
\end{figure}


Amulets of symbols were divided into two graphs, as those depicting \textit{udjat} were one of the most prevalent motifs. Quadruple \textit{udjat} amulets were also included in the same chart as the singular \textit{udjat} (Figure \ref{fig:phase-udjat}). Proportionally, \textit{udjat} amulets were mainly popular during the 25th Dynasty and Early Napatan phases, especially in elite and non-elite contexts.


The percentage from the pre-25th Dynasty phase in royal burials is very high because the only other amulets, apart from the combined Pataikos/hawk-headed deity, were of the \textit{udjat}. From the Early Napatan royal contexts, the quantity of \textit{udjat} amulets decreases, reaching the lowest proportion in the Middle Napatan phase. However, it spikes again in Late Napatan tombs, becoming proportionally higher than on the 25th Dynasty. Elite \textit{udjat} amulets display a clear decrease in use after the peak during the 25th Dynasty. Although they were still very common, making up 51\% in the 25th Dynasty, 32\% in Early Napatan, 30\% in Middle Napatan and 27\% in Late Napatan. Similarly, \textit{udjat} amulets were more predominant in non-elite burials from the 25th Dynasty and Early Napatan (45\% and 35\%, respectively). Pre-25th Dynasty and Middle Napatan \textit{udjat} amulets were rarer during the pre-25th Dynasty and Middle Napatan, 9\% and 12\%, respectively. The quadruple \textit{udjat} was only found in non-elite contexts, although they were rare---2\% of the total number of amulets from the 25th Dynasty and 5\% from the Early Napatan phase.


\begin{figure}[H]
	\centering
	\includegraphics[width=1\textwidth]{phase-udjat}
	\caption{Proportion (in percentage) of \textit{udjat} and quadruple \textit{udjat} amulets by social group and chronological phase, relative to the total number of amulets.}
	\label{fig:phase-udjat}
\end{figure}


Regarding the additional symbols amulets, they were separated into subcategories (Figure \ref{fig:phase-symbols}). These include winged motifs (winged griffins, scarabs, snakes and \textit{uraeus}); hieroglyphic signs; royal symbols (crook, was scepter, whip, \textit{uraeus}, double \textit{uraeus}, \textit{hmhm} crown and double feather); symbols with animals heads (\textit{uraei} with lion heads and ram-headed sphinxes); symbols related to the sun (\textit{akhet}, sun boat, sun disc, and horned sun disc); moon-related symbols (horned lunar disc and lunar crescent); common symbols which are very recurrent in amulets across Egypt and Nubia (\textit{menat}, heart, \textit{djed}, and Isis knot); single and double \textit{ba} birds and, finally, sphinxes.


All royal phases contained a certain amount of these amulets, with the most popular being common symbols during the 25th Dynasty, Early and Middle Napatan phases. The \textit{menat} was particularly prominent, as mentioned in Chapter 4. The other subcategories were rarer, although there is a slight increase in royal symbols, \textit{ba} birds and winged motifs in the Late Napatan burials. 25th Dynasty tombs also contained the most diverse assemblage, with seven out of nine subcategories represented---sun and moon-related motifs were absent. The majority of amulets represented winged motifs and common symbols. On the other hand, Early Napatan symbol amulets depicted only common and royal symbols, while those from Middle Napatan included mostly common symbols, but also hieroglyphic signs, \textit{ba} birds, and royal and sun-related motifs.


Elite contexts follow the same pattern, with common motifs more frequent in all phases, although they represented less than 40\% of the total amulets in this group. In fact, during the 25th Dynasty and Late Napatan, they were the only symbol amulets apart from the \textit{udjat}. Unlike royal contexts, only one \textit{menat} was found in elite burials. Most of the common motifs were composed of \textit{uraei}, hearts and \textit{djed}. Early Napatan assemblages contained the most diverse set of subcategories, encompassing common symbols, winged motifs, \textit{ba} birds, royal symbols, hieroglyphic signs and sun-related symbols. Overall, elite contexts were less diverse regarding these amulet types.


Apart from \textit{udjat} and quadruple \textit{udjat}, symbol amulets patterns of the non-elite are completely different from the other social groups. Firstly, they were only found in burials from the 25th Dynasty to Early Napatan phases. Secondly, they represented less than 10\% of the total amulets in each phase, and finally, royal motifs were the most frequent during the 25th Dynasty, while in the Early Napatan period, royal motifs and hieroglyphic signs were slightly more recurring than other motifs. The other subcategories were present but in significantly lower quantities, some with a unique example. Although royal symbols and hieroglyphic signs were proportionally more common, they were still very rare.


\begin{figure}[H]
	\centering
	\includegraphics[width=1\textwidth]{phase-symbols}
	\caption{Proportion (in percentage) of symbol amulets, divided in subcategories, by social group and chronological phase, relative to the total number of amulets.}
	\label{fig:phase-symbols}
\end{figure}


Animal amulets were divided into six subcategories (Figure \ref{fig:phase-animals}):

\begin{itemize}
	\item Animals associated with common Egyptian deities. In the corpus, these include cats (Bastet), jackals (Anubis), ibises and apes (Thoth), crocodiles (Sobek), hippopotamus (Taweret), lions (Maahes, Sekhmet and Tefnut), rams (Amun), snakes (Nehebkau), and falcons (Horus). Although, it is also important to consider that many of these animals could be found in nature and, therefore, were not necessarily associated with an Egyptian deity.
	\item Common animals, which are not associated with popular deities from the Egyptian pantheon, such as ducks, fishes, flies, frogs, giraffes, geese, hares, hedgehogs, lizards, monkeys, pigs, scorpions, sows, turtles, and unidentified birds.
	\item Double animals include double frogs and double ram motifs.
	\item Quadruple animals encompasses amulets with four apes and four-headed rams.
	\item Hybrid animals include amulets that combine elements from different animals, such as hawk-headed crocodiles, ram-headed scarabs, and lion-headed flies.
	\item Parts of animals, such as feathers and vulture wing motifs.
\end{itemize}


With a few exceptions, animals associated with Egyptian deities were the majority in all social groups. In royal contexts, only three subcategories were present in three phases, as no animal amulets were found in the pre-25th Dynasty and Late Napatan phases. 25th Dynasty and Early Napatan amulets were not very diverse, including only motifs of animals associated with Egyptian deities and common animals. Although during the Middle Napatan phase, animals related to Egyptian deities were still the most popular, common animals were replaced by quadruple animals, more specifically, a single amulet of a four-headed ram.


Although animal amulets were more prevalent in elite contexts, they were limited to those linked to Egyptian deities and common animals, mirroring the pattern of early royal assemblages. However, in 25th Dynasty burials, amulets of common animals were more prevalent than those linked to deities. While in Early and Middle Napatan phases, animals associated with deities were more frequent than common animals, Late Napatan contexts contained only animals linked to deities. There is a marked shift in the pattern for non-elite assemblages after the pre-25th Dynasty period, when only animals linked to Egyptian deities were found. Especially during the 25th Dynasty and Early Napatan phases, not only do common animals exceed deity-associated ones, but more subcategories are present. These include hybrid animals, and double and quadruple animals. Additionally, Early Napatan assemblages also contained amulets depicting parts of animals. Common animal amulets were also more popular than those linked to deities during the Middle Napatan phase.


\begin{figure}[H]
	\centering
	\includegraphics[width=1\textwidth]{phase-animals}
	\caption{Proportion (in percentage) of symbol amulets, divided in subcategories, by social group and chronological phase, relative to the total number of amulets.}
	\label{fig:phase-animals}
\end{figure}


\begin{figure}[H]
	\centering
	\includegraphics[width=1\textwidth]{phase-nature}
	\caption{Proportion (in percentage) of nature amulets, divided in subcategories, by social group and chronological phase, relative to the total number of amulets.}
	\label{fig:phase-nature}
\end{figure}


As seen before, amulets of natural elements were very rare in all social groups and phases, less than 10\% of the total (Figure \ref{fig:phase-nature}). Nature amulets were divided into two subcategories: common and symbolic plants. The first includes motifs of double leafs, flowers and trees, while the latter contains papyrus columns, lotus flowers and pomegranates, which incorporated symbolic meanings in Egyptian religion. In royal contexts, symbolic plants, mostly the papyrus column, were found in assemblages from the 25th Dynasty to Middle Napatan phases, and only in the latter there were some amulets of common plants. The highest distribution of nature amulets comes from elite contexts, also spanning the same phases as the royal ones. However, only symbolic plants were present, encompassing 3\% of all amulets in the 25th Dynasty, 9\% in Early Napatan and 6\% in Middle Napatan. In non-elite graves, those from the 25th Dynasty and Early Napatan phases contained mostly symbolic plants (1.2\% and 2.5\%, respectively), and single amulets of common plants. Conversely, during the Middle Napatan phase, they were only amulets of common plants, albeit they accounted for less than 1\% of the total.


\begin{figure}[H]
	\centering
	\includegraphics[width=1\textwidth]{phase-human}
	\caption{Proportion (in percentage) of human amulets, divided in subcategories, by social group and chronological phase, relative to the total number of amulets.}
	\label{fig:phase-human}
\end{figure}


Amulets depicting humans were exclusive to non-royal contexts (Figure \ref{fig:phase-human}). These motifs were divided into human figures (boys, twin boys, unidentified figures, naked females, man on horse and male with double crown) and parts of the body (eye, double eye, face and hands). Only parts of the body were included in very small numbers in elite burials of the 25th Dynasty and Early Napatan phases (1\% and 0.5\%, respectively). Meanwhile, contemporary non-elite assemblages contained human figures and parts of the body, with the latter being the most prevalent. 25th Dynasty human figures totaled 2\% and body parts, 2.6\%. In Early Napatan tombs, the percentages were 0.5\% for human figures and 2\% for parts of the body.  In contrast, in Middle Napatan contexts, only parts of the body were found in high quantities, making up more than 40\% of the total amulet count.


Finally, amulets representing objects and miscellaneous motifs were separated into six subcategories (Figure \ref{fig:phase-objs-misc}). For objects, they include professional objects (axe heads, dovetails, sistrums and harpoons); vessels (vases, situlas and jars); and funerary objects (headrests and offering tables). The miscellaneous category includes amulet cases, written motifs (such as cartouches), and non-figurative motifs (gold nuggets, pillars, tablets, an incised cylinder and pebbles). All these amulets were extremely rare, with the highest distribution barely reaching 4\% of the total. In the vast majority of contexts, these motifs encompass only a single amulet.


Royal amulets included the non-figurative inscribed gold nugget from a pre-25th Dynasty burial, and, from 25th Dynasty contexts, a single instance of another non-figurative amulet, one vessel and one funerary object. From Early Napatan, there was a single funerary object, a headrest amulet. Elite assemblages contained a single written motif amulet from the 25th Dynasty; one non-figurative motif, one funerary object and one written amulet from the Early Napatan; and one funerary object from Middle Napatan. In contrast, non-elite burials included all subcategories and most comprised of more than one amulet. 25th Dynasty amulets contained the same motifs as the previous groups, with the addition of professional objects and amulet cases. In fact, the former are the most frequent during the 25th Dynasty and Early Napatan phases, although still very rare---0.4\% and 0.8\% of all amulets, respectively. Conversely, Middle Napatan tombs contained the highest amount of these subcategories, but only included non-figurative motifs (2\%) and amulet cases (1.5\%).


\begin{figure}[H]
	\centering
	\includegraphics[width=1\textwidth]{phase-objs-misc}
	\caption{Proportion (in percentage) of object and miscellaneous amulets, divided in subcategories, by social group and chronological phase, relative to the total number of amulets.}
	\label{fig:phase-objs-misc}
\end{figure}


\begin{figure}[H]
	\centering
	\includegraphics[width=1\textwidth]{phase-2deities-1}
	\vspace{10pt}
	\includegraphics[width=1\textwidth]{phase-2deities-2}
	\caption{Distribution of deity amulets with two or more motifs by social group and chronological phase.}
	\label{fig:phase-2deities}
\end{figure}


As seen in previous chapters, many amulets combine more than one motif. However, here I want to explore the occurrence of amulets specifically depicting more than one deity. These amulets were only found in royal and non-elite burials, spanning from the pre-25th Dynasty (with the sole amulet being of a Pataikos and a hawk-headed deity) to Early Napatan phases (Figure \ref{fig:phase-2deities}). A recurrent type of amulet are back-to-back deities, especially those with either winged anthropomorphic or lion-headed goddesses. In royal contexts, they appear in a number of amulets from queen tombs of the 25th Dynasty. While only a single human-headed winged goddess is paired with a Pataikos, those with a lion head appear in six amulets, with different primary motifs. One of them is on the reverse of an amulet depicting the triad of Amun, Isis and Horus; another of Isis nursing Horus; while the other four are slightly related: two are of a Pataikos, another of a ram-headed dwarf and the last of a winged Pataikos. Other amulets with two deities from the 25th Dynasty include Amun with Ra, two amulets of the Pataikos with Shu on the back, and a Bes paired with Amun, Isis and Horus.


The same pattern can be observed in non-elite contexts. During the 25th Dynasty, Pataikos amulets were also the most common to be paired with other deities. In particular, they appear together with winged lion-headed goddesses, hawk-headed deities and the god Bes (two instances each). Other amulets to pair two deities include one with a hawk head and another with a crocodile head, as well as a winged goddess on the back of a Bes amulet. For burials spanning the 25th Dynasty and Early Napatan periods, three Pataikos amulets included a leonine winged goddess on the reverse and one were of two Pataikos back-to-back. The pairings Amun/Mut and Sekhmet/Nefertum were found in two amulets each. Other pairings include Khonsu and Thoth, Onuris and Tefnut, Ra-Horakhty and Nefertum, a winged goddess with a child, and, finally, Bes with the Theban triad (Amun, Mut and Khonsu), all with one instance each.


As these types of amulets were not abundant, and the ones with images available do not appear to have been made from the same mold, it is likely that they were not mass produced and were not available for everyone. However, the Pataikos seems to have been a particular popular deity to include another, especially on its back.


\subsection{Amulet materials}


\begin{figure}[H]
	\centering
	\includegraphics[width=1\textwidth]{phase-faience}
	\caption{Percentage of faience amulets by social group and chronological phase.}
	\label{fig:phase-faience}
\end{figure}


Faience is the most common material for amulets across most phases and social groups. Figure \ref{fig:phase-faience} illustrates the distribution of faience amulets in percentage according to social group and phase. In royal burials, faience amulets started to appear during the 25th Dynasty, as all amulets from the pre-25th Dynasty phase were made of gold. The frequency of faience increased from the 25th Dynasty to the Early Napatan, going from 51\% to 76\%. Afterwards, their prevalence declines, reaching 55\% in Middle Napatan and 6\% in Late Napatan (only one amulet out of the total 17). The opposite trend is observed with elite faience amulets. They were at an all time high during the 25th Dynasty (78\%), before collapsing to 31\% across Middle Napatan tombs, and then rising slightly in the Late Napatan phase (43\%). In contrast to the previous groups, in non-elite burials, faience first appears in pre-25th Dynasty amulets, with 45\% of the total. This number increases considerably during the 25th Dynasty (86\%), before declining very slightly in Early Napatan contexts (82\%), and reaching its highest proportion during the Middle Napatan phase (91\%).


As the faience workshop at Sanam Temple is considered to have been exclusively used to manufacture royal objects \parencite{howleySanam2026}, and given the high number of faience amulets in non-elite contexts, workshops must have existed to provide this strata of the population with these amulets and other faience goods. Additionally, workshops dated to the New Kingdom occupation could have been continuously used or reconstructed in subsequent periods. Furthermore, some faience amulets could also have been imported, either the material or the finished object.


\begin{figure}[H]
	\centering
	\includegraphics[width=1\textwidth]{phase-materials}
	\caption{Percentage of local and imported materials by social group and chronological phase.}
	\label{fig:phase-materials}
\end{figure}


The remaining materials were divided based on the location of the sources. Even without accounting for faience, the majority of amulets were made from materials that could be supplied from the Egyptian Eastern Desert and Lower Nubia (Figure \ref{fig:phase-materials}). More specifically, all royal amulets from the pre-25th Dynasty were made of gold, a local commodity, which could be explained by the collapse of the New Kingdom occupation. During the 25th Dynasty, most of the materials came either from the Eastern Desert, Lower Nubia or Asia, with single instances of agate from the Nile Valley and ivory from Sub-Saharan Africa. In contrast, the subsequent phases included materials solely sourced at the Eastern Desert and Lower Nubia, with a higher proportion for the latter, reaching 82\% in Late Napatan. However, these included few amulets in absolute numbers: 14 made of gold, one of quartz and one of serpentinite.


Elite contexts present a different pattern. For instance, in the 25th Dynasty, most materials (apart from faience) were sourced from Asia (10\%), followed by Lower Nubia (6\%) and Eastern Desert (5\%). During the Early Napatan phase, the majority of materials came from the Lower Nubia (35\%), succeeded by Eastern Desert (27\%), Asia (4\%) and one instance from the Nile Valley (sandstone). Middle to Late Napatan materials were mostly sourced from the Eastern Desert (27\% and 53\%, respectively) and Lower Nubia (12\% and 13\%) with a few instances coming from the Nile Valley (sandstone and agate).


For non-elite amulets, the source of materials was much more diverse following the pre-25th Dynasty phase, when, apart from faience, there were only materials from Lower Nubia (gold and carnelian), making up 55\% of amulets. 25th Dynasty and Early Napatan amulet assemblages include materials from Lower Nubia (5\% and 6\%, respectively), Eastern Desert (3\% and 5\%), Asia (2.3\% each), Nile Valley (0.8\% and 0.9\%), Eastern Sudan (0.1\% and 0.3\%), Sub-Saharan Africa (0.1\% and 0.2\%) and, finally, glass, which was classified as "Manufacturing", as it could be produced in a similar way to faience (1\% and 1.1\%). Additionally, there was one amulet from the 25th Dynasty made of mother-of-pearl, which could be found in the Red Sea. In contrast, Middle Napatan materials were limited to Eastern Desert and Lower Nubia (0.5\% and 8\%, respectively).  Although non-elite contexts included a greater diversity of materials coming from different sources, they were not the norm, as the great majority of amulets were made of faience.


\section{Regional variations}


In this section, amulets will be analyzed based on the geographical location of the sites in order to build regional objectscapes. Elite and non-elite burials were distributed across different locations in the Nubian territory, from Lower Nubia to the Meroe region, while royal tombs were only found in the 4th Cataract region, where the capital was located.


\subsection{Amulet distribution}


\begin{figure}[H]
	\centering
	\includegraphics[width=1\textwidth]{region-distro}
	\caption{Distribution (in percentage) of tombs containing amulets by social group and region.}
	\label{fig:region-distro}
\end{figure}


The frequency of amulets by region is not much different from the chronological distribution (Figure \ref{fig:region-distro}). For all royal tombs at the 4th Cataract, less than half of the burials contained amulets (44\%), while the amount was much lower in elite tombs from both the 4th Cataract and Meroe region, where only 9\% and 13\% of burials were furnished with these objects. Non-elite contexts, on the other hand, vary depending on the region, with tombs further south increasingly containing more amulets. More specifically, graves with amulets totaled 9\% in Lower Nubia, 13\% in North Upper Nubia and 23\% in the 4th Cataract and Meroe regions.


\subsection{Amulet types}

Regional amulet types mirror chronological patterns, with symbols, deities, and animals dominating---with the exception of Lower Nubia non-elite burials, where deities slightly outnumber symbols as the most frequent type (Figure \ref{fig:region-types}). However, across regions, assemblages show greater diversity than chronologically. In 4th Cataract royal burials, the majority of amulet types were present, with the exception of human motifs. Although nature, miscellaneous and object amulets were very rare. Elite contexts also varied. In the 4th Cataract, the only absent type is objects, whereas Meroe region tombs did not contain any human motifs. Lower Nubia and Meroe non-elite amulets were the most diverse, with all types present, although some were very rare, mostly nature and object motifs. North Upper Nubia was the least diverse, comprising of only three types: symbols, deities and animals. At the 4th Cataract, amulet types did not include miscellaneous motifs.


\begin{figure}[H]
	\centering
	\includegraphics[width=1\textwidth]{region-types}
	\caption{Frequency (in percentage) of amulet categories by social group and region.}
	\label{fig:region-types}
\end{figure}


\subsection{Amulet motifs}


Following the division of motifs established in the previous section, deity amulets from the Egyptian pantheon were also predominant in all regions (Figure \ref{fig:region-deities}). 4th Cataract royal amulets comprised 13\% of Egyptian deities, against 5\% of local deities or adaptations, the highest frequency across all social groups. These motifs were completely absent from elite tombs from the same location, while Egyptian deities totaled 30\%. Meroe elite amulets were mostly of Egyptian deities (14\%), with only 1\% of local deities/adaptations. Non-elite amulets follow the same pattern, with locations further south containing more local deities/adaptations. For instance, they were absent from Lower and North Upper Nubian burials, where Egyptian deities made up 30\% and 18\% of all amulets, respectively. In the 4th Cataract, local deities/adaptations do appear, but very rarely, only 0.1\% of the total, against 33\% of Egyptian deities. At the Meroe region, local deities/adaptations were a little more common (0.6\%), however Egyptian deities still dominated (22\%).


\begin{figure}[H]
	\centering
	\includegraphics[width=1\textwidth]{region-deities}
	\caption{Proportion (in percentage) of amulets of Egyptian deities and local adaptations by social group and region, relative to the total number of amulets.}
	\label{fig:region-deities}
\end{figure}


\begin{figure}[H]
	\centering
	\includegraphics[width=1\textwidth]{region-udjat}
	\caption{Proportion (in percentage) of \textit{udjat} and quadruple \textit{udjat} amulets by social group and region, relative to the total number of amulets.}
	\label{fig:region-udjat}
\end{figure}


Regarding \textit{udjat} amulets, they continue to be the most frequent in each location (Figure \ref{fig:region-udjat}). Royal amulets from the 4th Cataract encompassed 32\% of \textit{udjat}. In elite contexts, they totaled 40\% at the 4th Cataract and 36\% at the Meroe cemeteries. At non-elite burials, their frequency varied significantly by region. At Lower Nubia, they comprised 23.5\% of all amulets, and at North Upper Nubia, 22.5\%. The distribution increases in southern sites, reaching 35\% in the 4th Cataract and half of all amulets in the Meroe region.


The quadruple \textit{udjat}, as seen previously, was exclusive to non-elite burials from the 25th Dynasty to Early Napatan phases. Regionally, they were found in all regions, however in very different numbers. At Lower Nubia, they totaled 4\% of all amulets, whereas, at North Upper Nubia, they were more common than the standard \textit{udjat}, comprising 37\% of the total. Their presence decreases in the other sites, with 3\% at the 4th Cataract and 0.1\% in the Meroe region, which included only two examples.


\begin{figure}[H]
	\centering
	\includegraphics[width=1\textwidth]{region-symbols}
	\caption{Proportion (in percentage) of symbol amulets, divided in subcategories, by social group and region, relative to the total number of amulets.}
	\label{fig:region-symbols}
\end{figure}


Figure \ref{fig:region-symbols} illustrates the distribution of the remaining symbol amulets. The subgroups of these amulets, although being the most popular type in virtually all regions and social groups, display an interesting variability in their frequency. In 4th Cataract royal burials, while common symbols were more prevalent (33\%), other seven subcategories were also present, with only moon-related motifs absent. The second most common subcategory are winged motifs (3\%), followed by royal symbols (2\%). The remaining motifs were very rare, some with only a single example. More specifically, hieroglyphic signs made up 0.5\% of the total; single/double \textit{ba} birds, 0.4\%; symbols with animal heads 0.3\%; and sun-related and sphinx motifs, 0.1\% each.


Of these motifs, 4th Cataract elite amulets included only common symbols (2\%), while the Meroe cemeteries showcase a greater variety. For instance, common symbols were the most frequent (11\%), followed by royal symbols (5\%), hieroglyphic signs (4\%), winged motifs and sun-related motifs (1\% each), and finally, \textit{ba} birds (0.7\%). Moon-related motifs were the only absent subcategory.


Apart from the \textit{udjat}, non-elite contexts included, proportionally, less symbol amulets than other social groups. Lower Nubia, for instance, only included royal symbols (3\%), while, North Upper Nubia contained common symbols (5\%), followed by winged and sun-related motifs (0.8\% each). The 4th Cataract was more diverse, with hieroglyphic signs being slightly more frequent (1.4\%), followed by royal symbols (1\%). The remaining subcategories were very rare. Common symbols made up only 0.4\%; winged motifs and symbols with animal heads, 0.3\% each; sun- and moon-related motifs, 0.1\% each; and sphinxes with only one instance, 0.04\% of the total. Although the majority were extremely rare, this was the only context to include all subcategories. Non-elite amulets in the Meroe region included motifs of royal symbols (5\%),  hieroglyphic signs and common symbols (0.1\% each), and winged motifs, symbols with animal heads, \textit{ba} birds, and sun-related motifs, all with a single example (0.05\% each).


In contrast, the distribution of animal amulets subcategories was more homogeneous (Figure \ref{fig:region-animals}). At the 4th Cataract royal sites, animals associated with deities were the most frequent, 5\% of the total amulets. Common animals were rarer, only 0.8\%, while a single quadruple animal was found (0.1\%). Elite amulets from the same region were composed solely of common animals (12\%) and those associated with deities (5\%). At the Meroe cemeteries, the distribution is reversed, with animals linked to Egyptian deities dominating at 11\%, followed by common animals at 4\%. No other animal subcategories were found at elite contexts.


Non-elite animal amulets were the most diverse at the southern sites. Lower Nubian motifs encompassed 14\% of animals associated with deities and 4\% of common animals. Similarly to elite contexts, the distribution is reversed at North Upper Nubia, with common animals as the most frequent (10\%), followed by those linked to deities (4\%). Amulets from the 4th Cataract included all five subcategories: common animals (10\%), those associated with deities (6\%), hybrid animals (0.5\%), double animals (0.3\%), and quadruple animals (0.1\%). Similarly, amulets from the Meroe region also included all subcategories, albeit in fewer amounts. Animals associated with Egyptian deities were slightly more frequent (5\%), followed by common animals (3\%), hybrid animals (1.4\%), and double and quadruple animals with only one amulet each (0.05\% of the total each).


\begin{figure}[H]
	\centering
	\includegraphics[width=1\textwidth]{region-animals}
	\caption{Proportion (in percentage) of animal amulets, divided in subcategories, by social group and region, relative to the total number of amulets.}
	\label{fig:region-animals}
\end{figure}


Nature and human amulets were combined in the same distribution graph, as they were not very numerous (Figure \ref{fig:region-nature-human}). Amulets representing nature elements were mostly composed of symbolic plants, specifically the papyrus column. In 4th Cataract royal burials, they totaled 2\% of the total amulets, whereas common plants were only 0.1\%, with a single example. Elite cemeteries in the 4th Cataract and Meroe region included only symbolic plants (7\% each). Non-elite nature amulets from Lower Nubia included one example of a symbolic plant (1\% of the total). At the 4th Cataract, they were the majority, at 3\%, with only one instance of a common plant (0.04\%). At the Meroe region, common plants were also composed of only one amulet (0.1\%), and symbolic plants made up 3\% of the total amulets. No nature amulets were found in North Upper Nubian non-elite burials.


Amulets depicting humans and parts of the body were only found in elite and non-elite contexts. 4th Cataract elite amulets included only parts of the body, 3\% of the total. At non-elite sites, Lower Nubia also contained parts of the body (7\%), while 4th Cataract and Meroe region burials were slightly more diverse. At the 4th Cataract, parts of the body comprised 2\% of amulets, whereas human figures were 0.2\%, with only five examples. They were higher in the Meroe cemeteries (4\%), but parts of the body were still more prevalent (5\%). Similarly to nature amulets, North Upper Nubia sites did not include any human motifs.


\begin{figure}[H]
	\centering
	\includegraphics[width=1\textwidth]{region-nature-human}
	\caption{Proportion (in percentage) of nature and human amulets, divided in subcategories, by social group and region, relative to the total number of amulets.}
	\label{fig:region-nature-human}
\end{figure}


Miscellaneous amulets and those depicting objects were the rarest type, with the highest distribution barely reaching 2.5\% of the total amulets (Figure \ref{fig:region-objs-misc}). Among object motifs, 4th Cataract royal amulets included funerary items (0.3\%, two examples) and vessels (0.1\%, one example). At elite sites, only those from the Meroe region featured these motifs, specifically a single funerary object (0.7\%). Lower Nubian non-elite amulets included one professional object (1\%). As with the previous motifs, amulets from the southernmost sites were the most diverse. At the 4th Cataract, they included professional objects (0.5\%, eleven examples), funerary items (0.2\%, five examples), and vessels (0.1\%, three examples). In the Meroe cemeteries, the distribution is similar. Professional objects totaled 0.5\%, (ten items); and funerary objects and vessels were represented by only one amulet each (0.05\%).


\begin{figure}[H]
	\centering
	\includegraphics[width=1\textwidth]{region-objs-misc}
	\caption{Proportion (in percentage) of object and miscellaneous amulets, divided in subcategories, by social group and region, relative to the total number of amulets.}
	\label{fig:region-objs-misc}
\end{figure}


The miscellaneous category include non-figurative amulets, such as the nuggets, written motifs and amulet cases. One cartouche was found in the 4th Cataract royal cemeteries, while the elite sites at the same location also included one cartouche (2.4\% of the total). At the Meroe region, there was one non-figurative amulet (0.7\%). 4th Cataract non-elite sites included four written motifs (0.2\%), while at the Meroe region contained ten non-figurative motifs (0.5\%) and four amulet cases (0.2\%).


The frequency of amulets combining more than one deity can also be analyzed through a regional lens, especially with comparisons between the non-elite sites (Figure \ref{fig:region-2deities}). For instance, 4th Cataract non-elite burials included the most amulets with these dual motifs, nine in total. These include amulets of Amun/Mut and Sekhmet/Nefertum (both with two amulets each), Khonsu/Thoth, Onuris/Tefnut, Pataikos/Pataikos, Ra-Horakhty/Nefertum, Bes/winged goddess (all with one instance each). North Upper Nubia amulets featured combinations of Bes with Amun, Mut and Khonsu, and a winged goddess with a child. The pair Pataikos/winged leonine goddess was only found in the Meroe region (five examples). In this locality, the Pataikos was also paired with hawk-headed deities (two examples) and the god Bes (two amulets). Additionally, a Bes amulet included a anthropomorphic winged goddess on the reverse, and there was an amulet featuring two deities, one with a hawk head and another with a crocodile head.


In contrast, it was at 4th Cataract royal tombs that the Pataikos was the most paired deity, with six amulets. One was with a winged goddess, two with the god Shu, one with a hawk-headed deity and two with winged leonine goddesses. Overall, the 4th Cataract yielded the most double sided deity amulets, both for royal and non-elite. The later encompassed a greater variety of deities, including some not present at all in royal contexts, such as Nefertum (with Sekhmet and Ra-Horakhty), Onuris with Tefnut, Kohnsu with Thoth, and Amun with Mut. The main difference between these two contexts is the number of Pataikos amulets with other deities, as they were most frequent in Meroe non-elite sites. Nevertheless, the majority of deity amulets depicted those associated with the Egyptian pantheon. However, they could be paired with unidentified deities that could potentially be local adaptations, such as the winged goddesses and hawk- and crocodile-headed deities.


\begin{figure}[H]
	\centering
	\includegraphics[width=1\textwidth]{region-2deities-1}
	\vspace{10pt}
	\includegraphics[width=1\textwidth]{region-2deities-2}
	\caption{Distribution of deity amulets with two or more motifs by social group and chronological phase.}
	\label{fig:region-2deities}
\end{figure}


\subsection{Amulet materials}


\begin{figure}[H]
	\centering
	\includegraphics[width=1\textwidth]{region-faience}
	\caption{Percentage of faience amulets by social group and region.}
	\label{fig:region-faience}
\end{figure}


Regionally, the frequency of faience amulets varies by social group (Figure \ref{fig:region-faience}). While 55\% of 4th Cataract royal amulets were made of faience, elite amulets from the same location reaches 88\%. This distribution decreases significantly in the Meroe region, where only 29\% of all amulets were made of faience. At non-elite contexts, the distribution is more uniform. At Lower Nubia, the percentage is 85\%, then it drops slightly at North Upper Nubia (78\%), and rises again at the 4th Cataract (84\%) and Meroe region (87\%).


Regarding the distribution of the remaining materials, they were divided as the previous section  (Figure \ref{fig:region-materials}). Apart from faience, 4th Cataract royal amulets came mostly from the Egyptian Eastern Desert (18\%) and Lower Nubia (17\%). Materials from Asia totaled 6\%, while those from the Nile Valley and Sub-Saharan Africa were represented by just one example (0.2\% each). Elite amulet materials from the 4th Cataract were less diverse, as the majority was sourced from Lower Nubia (5\%), Eastern Desert and Asia (2\% each). At the Meroe cemeteries, where a higher number of amulets were not made of faience, materials from Lower Nubia made up 31\% of the total, followed by those from the Eastern Desert (29\%), with fewer originating in Asia (6\%) and the Nile Valley (2\%). 


Non-elite materials from Lower Nubia follow the same pattern as previous contexts, coming mainly from Lower Nubia (8\%), Eastern Desert (4\%) and Asia (3\%). At North Upper Nubia, materials were sourced from the Eastern Desert (12\%), Lower Nubia (6\%), and Nile Valley (2\%). In contrast, at the 4th Cataract sites, there is a slight increase in the diversity of materials, although in lesser numbers than previous regions. For instance, Lower Nubian materials accounted for less than 5\% of all amulets, followed by those from the Eastern Desert  (4\%). Materials from Asia, Upper Egypt, Nile Valley, and Eastern Sudan were even rarer (2\%, 1.5\%, 0.8\%, and 0.3\%, respectively). Additionally, 1.3\% of amulets were manufactured, which include 30 glass amulets. Finally, Meroe amulets were also diverse in terms of materials, albeit some included very few examples. The majority came from Lower Nubia (6\%), followed by Asia (2.3\%), Eastern Desert (2\%), Nile Valley (0.6\%), Sub-Saharan Africa (five examples of ivory, 0.3\%), glass manufacturing (two examples, 0.1\%), and the Red Sea (one mother-of-pearl amulet, 0.05\%).


\begin{figure}[H]
	\centering
	\includegraphics[width=1\textwidth]{region-materials}
	\caption{Percentage of local and imported materials by social group and region.}
	\label{fig:region-materials}
\end{figure}


\section{Conclusion}


This chapter compiled the data presented in the previous chapters, in order to compare it through three main focal points: social groups, chronological phase and geographical location. Through this analysis, certain patterns and deviations were noted.


Chronologically, the 25th Dynasty-Early Napatan period emerges as the most diverse for amulet assemblages, especially in non-elite groups. The most consistent distinction was between royal/elite and non-elite groups. Royal and elite amulets were dominated by symbol motifs, while non-elite contexts displayed a greater variety of types, with deity, animal and human amulets surpassing symbol motifs in certain phases. Furthermore, certain subcategories were exclusive to non-elite groups, such as the quadruple \textit{udjat}, moon-related motifs, amulet cases, professional objects, human figures (most concentrated in the 25th Dynasty-Early Napatan period). Although symbol amulets were not very numerous, it is significant that they include mostly royal symbols and hieroglyphic signs, in contrast to royal and elite groups, which were dominated by common symbols. Furthermore, non-elite groups included the least amount of local adaptations/deities. Amulets pairing two deities were limited to royal and non-elite groups, and were only included in graves from the 25th Dynasty-Early Napatan period. The majority of motifs were very rare, but the Pataikos with a winged goddess the most prevalent in both groups. Faience was the dominant material, especially for the non-elite. An interesting exception were royal amulets from the Late Napatan phase. The 25th Dynasty-Early Napatan period was the most materially diverse for all groups, especially for the non-elite. In later phases and in all groups, sources were limited to Lower Nubia and the Eastern Desert, while pre-25th Dynasty contexts used only faience and Lower Nubian materials.


Regionally, amulets from the 4th Cataract and Meroe region were the most varied, particularly in non-elite groups. Symbols amulets become more prevalent in non-elite groups, although deities still outnumber them at Lower Nubia and the 4th Cataract. Non-elite sites continue to be the most diverse, with the exception of those from North Upper Nubia. Some subcategories were restricted to the 4th Cataract and Meroe regions--oftentimes only in non-elite contexts--, such as local adaptation/deities, human figures, quadruple and double animals, hybrid animals, common plants, funerary objects, vessels, and non-figurative and written motifs. On the other hand, amulets such as the quadruple \textit{udjat} were absent from non-elite sites at the Meroe cemeteries. Although common symbols were prevalent in royal and elite sites from the 4th Cataract and Meroe regions, their higher frequency in non-elite was limited to North Upper Nubia. Common animals were more numerous than those linked to Egyptian deities at 4th Cataract elite and non-elite contexts, as well as non-elite North Upper Nubian burials. Apart from the royal sites at the 4th Cataract, amulet with two deities were present at non-elite sites in North Upper Nubia, 4th Cataract and Meroe region, with the later being the only one yielding the Pataikos with a winged goddess. Faience was, once again, the most common material for all regions, with the exception of Meroe elite burials. The majority of materials were sourced from the Eastern Desert, Lower Nubia and Asia. Nevertheless, 4th Cataract and Meroe region non-elite contexts were the most diverse, with materials coming from six different territories, and also included glass manufacture.